\documentclass{article}
\usepackage{amsmath, amssymb, amsthm}
\usepackage{hyperref}
\usepackage{geometry}
\geometry{margin=1in}

\title{Erd\H{o}s Problem \#145}
\date{Last updated: 2025-08-31}
\author{Source: erdosproblems.com}

\begin{document}
\maketitle

\noindent\textbf{Status:} OPEN \quad \textbf{No prize}

\noindent\textbf{Tags:} number theory

\medskip

\noindent\textbf{Problem Statement:}

Let $s_1&#60;s_2&#60;\cdots$ be the sequence of squarefree numbers. Is it true that, for any $\alpha \geq 0$,\[\lim_{x\to \infty}\frac{1}{x}\sum_{s_n\leq x}(s_{n+1}-s_n)^\alpha\]exists?




Erd\H{o}s \cite{Er51} proved this for all $0\leq \alpha \leq 2$, and Hooley \cite{Ho73} extended this to all $\alpha \leq 3$. Greaves, Harman, and Huxley showed (in Chapter 11 of \cite{GHH97}) that this is true for $\alpha \leq 11/3$. Chan \cite{Ch23c} has extended this to $\alpha \leq 3.75$. Granville \cite{Gr98} proved that this follows (for all $\alpha \geq 0$) from the ABC conjecture . See also [208] .




References

[Ch23c] Chan, Tsz Ho, On moments of gaps between consecutive square-free numbers . Mosc. J. Comb. Number Theory (2023), 287--295.

[Er51] Erd\H{o}s, P., Some problems and results in elementary number theory . Publ. Math. Debrecen (1951), 103-109.

[GHH97] Greaves, G. R. H. and Harman, G. and Huxley, M. N., Sieve Methods, Exponential Sums, and their Applications in Number Theory . (1997).

[Gr98] Granville, Andrew, {$ABC$} allows us to count squarefrees . Internat. Math. Res. Notices (1998), 991--1009.

[Ho73] Hooley, Christopher, On the intervals between consecutive terms of sequences . (1973), 129-140.

\noindent\textbf{Related OEIS sequences:} A005117


\bigskip
\noindent\small{Source: \url{https://www.erdosproblems.com/145}}
\end{document}
